% !PRINTED none

\begin{center}
\gesperrt{Dreizehnter Abschnitt.}

\medskip
\gesperrt{Allgemeine Theorie der Gleichung f\"unften Grades.}

\medskip
\rule{0.3\linewidth}{0.4pt}

\medskip
\S.\,106.

\gesperrt{Fragestellung.}
\end{center}

Wir haben im \S.\,44 gesehen, dass es \"uber die Frage nach
der \gesperrt{Galois'schen} Gruppe einer algebraischen Gleichung noch
ein weiter gehendes Problem giebt, n\"amlich die Frage nach der
linearen Substitutionsgruppe von m\"oglichst geringer Dimensionenzahl
, auf deren Formenproblem die gegebene Gleichung zur\"uck-
f\"uhrbar ist. Bei den metacyklischen Gleichungen, insbesondere
also auch bei den allgemeinen Gleichungen 3\textsuperscript{ten} und 4\textsuperscript{ten} Grades,
ist diese Dimensionenzahl gleich 1, wodurch eben ausgedr\"uckt
ist, dass diese Gleichungen durch Radicale l\"osbar sind. Die
zun\"achst zu untersuchenden Gleichungen sind dann die vom
5\textsuperscript{ten} Grade, und es wird sich zeigen, dass die L\"osung der allgemeinen
 Gleichung 5\textsuperscript{ten} Grades auf eine bin\"are lineare Substitutionsgruppe
, n\"amlich auf das Ikosa\"ederproblem f\"uhrt.

Damit im Zusammenhange steht aber noch eine andere Frage.
Die allgemeine Gleichung 5\textsuperscript{ten} Grades enth\"alt f\"unf Co\"efficienten
, die als unabh\"angige Variable betrachtet werden k\"onnen,
und folglich ist die Wurzel einer solchen Gleichung eine algebraische
 Function von f\"unf Variablen. Nun kann man aber
schon durch die einfachen linearen Substitutionen die Zahl dieser
Variablen vermindern. Durch \gesperrt{Tschirnhausen}-Transformation,
z.\,B. auf die \gesperrt{Jerrard'sche} Form, kann man die Gleichung sogar
nur von einem variablen Co\"efficienten (einem \gesperrt{Parameter}) ab-
h\"angig machen, und man kann also, durch Vermittelung von
Gleichungen, die den 5\textsuperscript{ten} Grad nicht erreichen, die Wurzel einer
allgemeinen Gleichung 5\textsuperscript{ten} Grades von einer algebraischen
Function von einer Variablen abh\"angig machen.
